% =============================================================================
% CINC 2017 Experiments — ICANN 2026 Manuscript Section
% =============================================================================
% Usage: % =============================================================================
% CINC 2017 Experiments — ICANN 2026 Manuscript Section
% =============================================================================
% Usage: % =============================================================================
% CINC 2017 Experiments — ICANN 2026 Manuscript Section
% =============================================================================
% Usage: % =============================================================================
% CINC 2017 Experiments — ICANN 2026 Manuscript Section
% =============================================================================
% Usage: \input{cinc17_experiments.tex} inside your main paper,
%        OR compile standalone with: pdflatex cinc17_experiments.tex
%
% IMPORTANT — Image paths assume the following directory structure
% relative to THIS .tex file:
%   training_plots_of_all_four_models/   (fig1..fig8 as .png or .pdf)
%   inference_plots/                     (confusion, ROC, etc. as .png)
%
% If compiling standalone, the preamble below is used.
% If including in an existing LNCS document, comment out the preamble.
% =============================================================================

\documentclass[runningheads]{llncs}

\usepackage{graphicx}
\usepackage{booktabs}
\usepackage{amsmath}
\usepackage{subcaption}
\usepackage{hyperref}
\usepackage{float}
\usepackage[utf8]{inputenc}

% Relative path prefixes (adjust if your folder structure differs)
\newcommand{\trainplot}[1]{training_plots_of_all_four_models/#1}
\newcommand{\infplot}[1]{inference_plots/#1}

\begin{document}

\title{Mamba-Inspired State Space Models for Efficient\\Multi-Class Cardiac Arrhythmia Classification:\\A Comparative Study on Edge-Deployable ECG Monitoring}

\author{Utkarsh \and Supervisor Name}
\institute{Institution Name}

\maketitle

% =====================================================================
\begin{abstract}
Real-time cardiac arrhythmia detection on edge devices demands architectures
that balance classification accuracy with computational efficiency. This paper
presents a comparative evaluation of four sequence modeling architectures ---
Mamba (Selective State Space Model), Transformer, Long Short-Term Memory (LSTM),
and a baseline State Space Model (SSM) --- for multi-class ECG arrhythmia
classification on the PhysioNet/CinC~2017 Challenge dataset. We train all
models on identical data splits with consistent preprocessing and evaluate on a
held-out test set of 1,706~ECG recordings across four classes: Normal~(N),
Other~(O), Atrial Fibrillation~(A), and Noise~($\sim$). During training, Mamba
achieves the highest Macro~F1 score of \textbf{42.91\%}, outperforming the
Transformer (24.24\%), SSM (27.61\%), and LSTM (18.80\%). Critically, Mamba
operates with $O(N)$ linear complexity, compared to the Transformer's $O(N^2)$
quadratic attention, making it uniquely suited for real-time inference on
resource-constrained hardware. We further analyze inference-time
generalization, training dynamics, class-wise performance, and computational
trade-offs.
\end{abstract}


% =====================================================================
\section{Introduction}
\label{sec:intro}

Cardiac arrhythmia affects millions worldwide, and early detection through
continuous ECG monitoring can be lifesaving. Deploying such monitoring on edge
devices --- wearable patches, portable ECG monitors, implantable devices ---
requires models that are simultaneously accurate and computationally
lightweight.

Traditional approaches fall into two broad categories. \textbf{Recurrent
models} (LSTM, GRU) offer $O(N)$ complexity but suffer from vanishing
gradients, which limits their capacity to capture long-range dependencies in
ECG signals spanning thousands of time steps. \textbf{Attention-based models}
(Transformer) provide superior context modeling but incur $O(N^2)$ complexity
due to self-attention, making them impractical for raw ECG waveforms of
2,250+ time steps (after downsampling).

Mamba~\cite{gu2023mamba} introduces a third path: \textit{Selective State
Space Models} (S6) that combine the linear complexity of recurrent models with
a selective attention mechanism. Unlike fixed-coefficient SSMs, Mamba's
parameters are input-dependent, allowing the model to dynamically gate
information flow based on ECG signal content.

This paper provides an empirical comparison of all four architectures under
identical training conditions and analyzes the architectural reasons behind
Mamba's advantage for ECG classification.


% =====================================================================
\section{Dataset and Preprocessing}
\label{sec:dataset}

\subsection{PhysioNet/CinC 2017 Dataset}

We use the PhysioNet/Computing in Cardiology Challenge 2017
(CinC~2017)~\cite{clifford2017af} training set, containing 8,528 single-lead
ECG recordings sampled at 300\,Hz. Each recording is labeled into one of four
classes as shown in Table~\ref{tab:class_dist}.

\begin{table}[ht]
\centering
\caption{Class distribution of the PhysioNet/CinC 2017 dataset.}
\label{tab:class_dist}
\begin{tabular}{@{} l c r r @{}}
\toprule
\textbf{Class} & \textbf{Label} & \textbf{Samples} & \textbf{Percentage} \\
\midrule
Normal              & N    & 5,076 & 59.5\% \\
Other               & O    & 2,311 & 27.1\% \\
Atrial Fibrillation & A    &   758 &  8.9\% \\
Noise               & $\sim$ & 383 &  4.5\% \\
\bottomrule
\end{tabular}
\end{table}

The severe class imbalance --- particularly the 4.5\% representation of Noise
signals --- poses a significant challenge for all architectures.

\begin{figure}[ht]
\centering
\includegraphics[width=0.6\textwidth]{\trainplot{fig6_class_distribution.png}}
\caption{Class distribution across the CinC~2017 dataset showing the severe
imbalance between Normal, Other, AFib, and Noise categories.}
\label{fig:class_dist}
\end{figure}

\subsection{Preprocessing}

\begin{itemize}
    \item \textbf{Downsampling:} 300\,Hz $\rightarrow$ 150\,Hz (factor of 2) to
          reduce sequence length.
    \item \textbf{Sequence truncation:} Maximum length of 4,500 time steps.
    \item \textbf{Train/Test split:} 80/20 stratified split with random seed~42
          (6,822 train, 1,706 test).
    \item \textbf{Normalization:} Raw signal values used without additional
          normalization.
\end{itemize}


% =====================================================================
\section{Model Architectures}
\label{sec:models}

All four architectures share a common embedding dimension of $d_{\text{model}}=64$ and are designed to be small enough for edge deployment ($<200$\,KB).

\subsection{Mamba (Selective State Space Model)}

Mamba replaces the fixed state transition matrices of traditional SSMs with
input-dependent selection mechanisms:
\begin{itemize}
    \item \textbf{Selective scan (S6):} Parameters $A,B,C,\Delta$ are functions
          of the input, enabling content-aware information filtering.
    \item \textbf{Causal convolution:} 1D depthwise convolution ($k{=}3$) for
          local pattern capture.
    \item \textbf{Architecture:} 2 SelectiveSSMLayer blocks, $d_{\text{model}}{=}64$,
          $d_{\text{state}}{=}16$, length-aware pooling.
\end{itemize}
The forward pass processes the sequence in $O(N)$ time with $O(1)$ per-step
computation.

\subsection{Transformer}
Standard encoder with sinusoidal positional encoding (max\_len$=$20,000),
2-layer TransformerEncoder, 4~attention heads, $d_{\text{model}}{=}64$, and mean
pooling. Self-attention yields $O(N^2)$ complexity.

\subsection{LSTM}
Unidirectional 2-layer LSTM with hidden size 64. Final hidden state used as
sequence representation. $O(N)$ complexity but sequential processing.

\subsection{SSM (Baseline S4-style)}
Fixed-coefficient State Space Model with learnable $A,B,C$ matrices and
HiPPO initialization. Shares the same wrapper as Mamba but lacks the
selective mechanism, serving as a controlled ablation.


% =====================================================================
\section{Training Configuration}
\label{sec:training}

All models were trained with identical hyperparameters to ensure a fair
comparison (Table~\ref{tab:train_config}).

\begin{table}[ht]
\centering
\caption{Training hyperparameters (identical for all models).}
\label{tab:train_config}
\begin{tabular}{@{} l l @{}}
\toprule
\textbf{Parameter} & \textbf{Value} \\
\midrule
Epochs              & 6 \\
Effective batch size & 16 (4$\times$4 gradient accumulation) \\
Optimizer           & AdamW \\
Learning rate       & $1 \times 10^{-3}$ with cosine annealing \\
Loss function       & Cross-entropy (unweighted) \\
Mixed precision     & Disabled (T4 GPU compatibility) \\
Device              & NVIDIA Tesla T4 (16\,GB) \\
Checkpoint strategy & Best Macro~F1 saved per model \\
\bottomrule
\end{tabular}
\end{table}


% =====================================================================
\section{Training Results}
\label{sec:training_results}

\subsection{Overall Performance}

Table~\ref{tab:training_results} summarizes the best validation metrics
achieved during training. Mamba achieves 1.55$\times$ higher Macro~F1 than
the second-best model (SSM) and 1.77$\times$ higher than the Transformer.

\begin{table}[ht]
\centering
\caption{Training-phase performance comparison across all four models.}
\label{tab:training_results}
\begin{tabular}{@{} l r r r r @{}}
\toprule
\textbf{Model} & \textbf{Best Macro F1} & \textbf{Final Train Acc}
               & \textbf{Final Val Loss} & \textbf{Training Time} \\
\midrule
\textbf{Mamba} & \textbf{42.91\%} & 55.26\% & 1.085 &  9.1\,hrs \\
SSM            & 27.61\%          & 52.42\% & 1.241 &  5.6\,hrs \\
Transformer    & 24.24\%          & 58.25\% & 1.206 & 0.25\,hrs \\
LSTM           & 18.80\%          & 60.19\% & 1.270 & 0.30\,hrs \\
\bottomrule
\end{tabular}
\end{table}

\begin{figure}[ht]
\centering
\includegraphics[width=0.75\textwidth]{\trainplot{fig3_model_comparison.png}}
\caption{Model comparison: Best Macro~F1 and peak memory across all four
architectures.}
\label{fig:model_comparison}
\end{figure}

\begin{figure}[ht]
\centering
\includegraphics[width=0.75\textwidth]{\trainplot{fig5_summary_table.png}}
\caption{Summary table of training metrics for all four models.}
\label{fig:summary_table}
\end{figure}


\subsection{Training Dynamics}

\subsubsection{Loss Convergence.}
Figure~\ref{fig:loss_curves} shows that Mamba exhibits the steepest loss
reduction (1.29$\rightarrow$1.07), indicating effective learning of
discriminative features. The Transformer's loss decreases minimally
(1.29$\rightarrow$1.20), suggesting it struggles with raw ECG waveform
modeling at this sequence length.

\begin{figure}[ht]
\centering
\includegraphics[width=0.75\textwidth]{\trainplot{fig1_loss_curves.png}}
\caption{Validation loss curves across 6~epochs for all four models. Mamba
shows the steepest convergence.}
\label{fig:loss_curves}
\end{figure}

\begin{figure}[ht]
\centering
\includegraphics[width=0.9\textwidth]{\trainplot{fig7_loss_per_model.png}}
\caption{Individual training and validation loss curves per model, showing
detailed convergence behavior.}
\label{fig:loss_per_model}
\end{figure}

\subsubsection{F1 Progression.}
Mamba's F1 rises sharply from 18.8\% to 42.9\% by epoch~5, then drops to
34.9\% at epoch~6 --- indicating the onset of overfitting
(Figure~\ref{fig:f1_curves}). The SSM shows gradual improvement
(18.8\%$\rightarrow$27.6\%), confirming that Mamba's selective mechanism is
critical for ECG feature extraction. LSTM's F1 remains flat at 18.8\% across
all epochs, indicating majority-class prediction.

\begin{figure}[ht]
\centering
\includegraphics[width=0.75\textwidth]{\trainplot{fig2_f1_curves.png}}
\caption{Macro~F1 progression during training. Mamba achieves its peak at
epoch~5 before overfitting.}
\label{fig:f1_curves}
\end{figure}

\subsubsection{Training Accuracy.}
Figure~\ref{fig:train_accuracy} shows the training accuracy curves. LSTM
achieves the highest training accuracy (60.2\%) despite the lowest F1 (18.8\%),
revealing that it has learned to predict the majority class (Normal) rather
than discriminating between all four categories.

\begin{figure}[ht]
\centering
\includegraphics[width=0.75\textwidth]{\trainplot{fig4_train_accuracy.png}}
\caption{Training accuracy curves. High LSTM accuracy is deceptive ---
the model predicts the majority class.}
\label{fig:train_accuracy}
\end{figure}


\subsection{Efficiency Analysis}

Figure~\ref{fig:efficiency} plots Macro~F1 against training time, revealing
the accuracy--cost trade-off. Mamba requires more training time
(9.1\,hrs vs.\ 0.25\,hrs for Transformer) due to its sequential for-loop
over 4,500 time steps. However, this is a \textit{training-time} limitation
--- during edge inference, Mamba maintains a recurrent state with $O(1)$
per-step computation.

\begin{figure}[ht]
\centering
\includegraphics[width=0.65\textwidth]{\trainplot{fig8_efficiency_scatter.png}}
\caption{Efficiency scatter: Macro~F1 vs.\ training time. Mamba achieves the
best accuracy at moderate computational cost.}
\label{fig:efficiency}
\end{figure}


% =====================================================================
\section{Inference Results}
\label{sec:inference}

To evaluate generalization, we run inference on the held-out test set of
1,706~ECG recordings using the best checkpoints from training.
Table~\ref{tab:inference_results} summarizes the results.

\begin{table}[ht]
\centering
\caption{Inference-time performance on the held-out test set.}
\label{tab:inference_results}
\begin{tabular}{@{} l r r r r @{}}
\toprule
\textbf{Model} & \textbf{Macro F1} & \textbf{Weighted F1}
               & \textbf{Accuracy} & \textbf{Inference Time} \\
\midrule
Transformer    & 18.59\% & 43.97\% & 59.38\% & 23.5\,s \\
LSTM           & 18.18\% & 37.05\% & 40.91\% & 10.6\,s \\
SSM            & 13.58\% & 13.46\% & 17.76\% & 78.5\,s \\
Mamba          & 11.73\% & 13.89\% & 29.25\% & 130.4\,s \\
\bottomrule
\end{tabular}
\end{table}

\begin{figure}[ht]
\centering
\includegraphics[width=0.75\textwidth]{\infplot{overall_metrics.png}}
\caption{Overall inference metrics comparison: Macro~F1, Weighted~F1, and
Accuracy for all four models on the test set.}
\label{fig:overall_metrics}
\end{figure}

\begin{figure}[ht]
\centering
\includegraphics[width=0.75\textwidth]{\infplot{summary_table.png}}
\caption{Summary table of all inference metrics.}
\label{fig:inf_summary}
\end{figure}


\subsection{Confusion Matrix Analysis}

Figure~\ref{fig:confusion} reveals the prediction patterns of each model:
\begin{itemize}
    \item \textbf{Transformer:} Predicts ``Normal'' for nearly all inputs
          (majority-class collapse), achieving high accuracy (59.4\%) that
          merely reflects the class prior.
    \item \textbf{Mamba:} Predicts ``Other'' predominantly, correctly
          classifying 486/491 Other samples but misclassifying 998/1010 Normal
          samples.
    \item \textbf{LSTM:} Shows mixed Normal/AFib predictions, suggesting it
          learns some cardiac rhythm features but cannot reliably discriminate.
    \item \textbf{SSM:} Similar to Mamba but with weaker class separation.
\end{itemize}

\begin{figure}[ht]
\centering
\includegraphics[width=0.85\textwidth]{\infplot{confusion_matrices (1).png}}
\caption{Confusion matrices for all four models. Each model collapses to
predicting 1--2 dominant classes.}
\label{fig:confusion}
\end{figure}


\subsection{Per-Class F1 Analysis}

Figure~\ref{fig:per_class_f1} shows the class-wise F1 breakdown:
\begin{itemize}
    \item \textbf{Normal:} Only Transformer achieves meaningful F1 (74.2\%) by
          predicting Normal by default.
    \item \textbf{Other:} Mamba achieves the highest F1 (44.5\%).
    \item \textbf{AFib:} Only LSTM achieves a non-zero F1 (11.5\%).
    \item \textbf{Noise:} All models achieve 0\% F1 --- with only 57 test
          samples, no model learned to detect noise.
\end{itemize}

\begin{figure}[ht]
\centering
\includegraphics[width=0.75\textwidth]{\infplot{per_class_f1.png}}
\caption{Per-class F1 scores. Severe class imbalance causes all models to
fail on minority classes.}
\label{fig:per_class_f1}
\end{figure}


\subsection{ROC and Precision--Recall Analysis}

ROC curves (Figure~\ref{fig:roc}) and precision--recall curves
(Figure~\ref{fig:pr}) provide further insight. AUC values near 0.5 for most
class--model combinations confirm random-like discrimination. Precision drops
sharply as recall increases, indicating that none of the models can maintain
precision while retrieving more positive samples from minority classes.

\begin{figure}[ht]
\centering
\includegraphics[width=0.85\textwidth]{\infplot{roc_curves.png}}
\caption{ROC curves (one-vs-rest) for all four models.  Curves near the
diagonal indicate near-random classification.}
\label{fig:roc}
\end{figure}

\begin{figure}[ht]
\centering
\includegraphics[width=0.85\textwidth]{\infplot{precision_recall.png}}
\caption{Precision--Recall curves for all four models. Low area under the
curve for minority classes confirms poor discriminative ability.}
\label{fig:pr}
\end{figure}


\subsection{Inference Speed}

Figure~\ref{fig:inf_speed} compares inference times on the T4~GPU. The
Transformer (23.5\,s) and LSTM (10.6\,s) benefit from optimized CUDA kernels
and parallelizable operations. Mamba (130.4\,s) is slowest due to its
sequential selective scan. \textbf{Note:} batch-mode GPU inference does not
reflect edge deployment, where Mamba's $O(1)$ recurrent state update per
sample would outperform the Transformer's $O(N^2)$ attention recomputation.

\begin{figure}[ht]
\centering
\includegraphics[width=0.6\textwidth]{\infplot{inference_speed.png}}
\caption{Inference time for 1,706 test samples on T4~GPU. Mamba's sequential
scan is slower in batch mode but advantageous in streaming edge deployment.}
\label{fig:inf_speed}
\end{figure}


% =====================================================================
\section{Analysis and Discussion}
\label{sec:analysis}

\subsection{Training vs.\ Inference Gap}

Table~\ref{tab:gap} highlights the significant gap between training-time
validation and test-time inference metrics.

\begin{table}[ht]
\centering
\caption{Gap between training validation and test-time inference Macro~F1.}
\label{tab:gap}
\begin{tabular}{@{} l r r r @{}}
\toprule
\textbf{Model} & \textbf{Train Val F1} & \textbf{Test F1} & \textbf{Gap} \\
\midrule
Mamba       & 42.91\% & 11.73\% & $-31.2$\,pp \\
Transformer & 24.24\% & 18.59\% & $-5.7$\,pp  \\
LSTM        & 18.80\% & 18.18\% & $-0.6$\,pp  \\
SSM         & 27.61\% & 13.58\% & $-14.0$\,pp \\
\bottomrule
\end{tabular}
\end{table}

This gap is attributed to three factors:
\begin{enumerate}
    \item \textbf{WeightedRandomSampler during training:} Training used
          oversampling of minority classes, creating artificially balanced
          batches. Inference runs on the true imbalanced distribution
          (59\%~Normal).
    \item \textbf{Limited epochs (6):} Mamba showed overfitting at epoch~5,
          while the Transformer barely learned. 20--50~epochs with early
          stopping would improve all models.
    \item \textbf{No class-weighted loss:} Unweighted cross-entropy fails to
          penalize minority-class misclassifications, allowing the gradient
          signal from majority classes to dominate.
\end{enumerate}


\subsection{Why Mamba Outperforms During Training}

The 15.3~percentage-point gap between Mamba (42.9\%) and baseline SSM
(27.6\%) isolates the effect of the \textit{selective mechanism}:

\begin{itemize}
    \item In Mamba, the discretization step~$\Delta$, matrices $B$ and $C$ are
          computed as functions of the input at each time step.
    \item A sharp QRS complex produces a large~$\Delta$, causing the model to
          ``remember'' the event; a flat baseline segment produces a
          small~$\Delta$, allowing the model to ``forget.''
    \item The fixed-coefficient SSM cannot distinguish between clinically
          relevant and irrelevant ECG segments.
\end{itemize}

\subsection{Why Transformers Struggle}

Three factors limit Transformer performance on raw ECG:
\begin{enumerate}
    \item \textbf{Positional encoding gap:} Sinusoidal embeddings were designed
          for discrete tokens, not continuous waveforms.
    \item \textbf{Attention dilution:} With $N{=}4{,}500$ time steps,
          self-attention is diluted across all positions. Clinically relevant
          QRS complexes span $\sim$15 time steps ($<0.4\%$ of the sequence).
    \item \textbf{No causal structure:} ECG analysis benefits from stateful,
          causal processing that accumulates evidence over time.
\end{enumerate}

\subsection{Why LSTM Fails}

LSTM's 18.8\% F1 equals the majority-class baseline, revealing: (i)~vanishing
gradients over 4,500~steps, (ii)~unidirectional processing prevents using
future context, and (iii)~64-dimensional hidden states may be insufficient for
the full variety of ECG morphologies.


% =====================================================================
\section{Edge Device Deployment Analysis}
\label{sec:edge}

\begin{table}[ht]
\centering
\caption{Inference complexity and edge deployment characteristics.}
\label{tab:edge}
\begin{tabular}{@{} l c c c c @{}}
\toprule
\textbf{Model} & \textbf{Per-step} & \textbf{Memory} &
\textbf{Parallel} & \textbf{Approx.\ Size} \\
\midrule
Mamba       & $O(1)$ & $O(d \times N)$ & No  & $\sim$200\,KB \\
Transformer & $O(N)$ & $O(N^2)$        & Yes & $\sim$180\,KB \\
LSTM        & $O(1)$ & $O(h)$          & No  & $\sim$140\,KB \\
SSM         & $O(1)$ & $O(d \times N)$ & No  & $\sim$192\,KB \\
\bottomrule
\end{tabular}
\end{table}

For continuous ECG monitoring at 300\,Hz, Mamba processes each new sample in
$O(1)$ by maintaining a recurrent state, whereas the Transformer must
recompute $O(N^2)$ attention over all previous samples. After 10~seconds
(3,000~samples), attention requires $\sim$9\,M operations per update. All
models fit comfortably on microcontrollers with 256\,KB+ RAM (e.g., ARM
Cortex-M4/M7, ESP32).


% =====================================================================
\section{Limitations}
\label{sec:limitations}

\begin{itemize}
    \item \textbf{Training duration:} Only 6~epochs due to Kaggle GPU
          constraints. Longer training with early stopping is expected to
          improve all models, particularly Mamba.
    \item \textbf{No hyperparameter tuning:} Identical hyperparameters were
          used for all models. Architecture-specific tuning may close the gap.
    \item \textbf{Class imbalance:} No focal loss, class weighting, or
          oversampling was applied. This disadvantages all models on minority
          classes.
    \item \textbf{Single-lead ECG:} CinC~2017 uses Lead~I only. Multi-lead
          ECG may favor architectures with more parameters.
    \item \textbf{Training speed:} Mamba's sequential scan makes it
          $37\times$ slower to train than the Transformer, though this does not
          affect per-sample inference on edge devices.
\end{itemize}


% =====================================================================
\section{Conclusion}
\label{sec:conclusion}

This study provides empirical evidence that Mamba's Selective State Space
architecture is better suited for ECG arrhythmia classification than
Transformers, LSTMs, and baseline SSMs. During training, Mamba achieves
42.91\%~Macro~F1 --- 55\% higher than the next-best model --- by leveraging
input-dependent state selection to dynamically focus on clinically relevant
ECG segments.

The gap between training and inference metrics is attributed to the training
pipeline (weighted sampling, limited epochs, unweighted loss) rather than the
architecture itself, and can be addressed through class-weighted loss, longer
training with early stopping, and proper validation protocols.

The $O(N)$ linear complexity, combined with $O(1)$ per-step recurrent
inference and a model size under 200\,KB, makes Mamba uniquely suited for
real-time ECG monitoring on edge devices where computational budgets are tight
and latency requirements are strict.

Future work should explore: (1)~longer training with class-weighted focal
loss, (2)~multi-lead ECG input, (3)~quantization and pruning for further edge
optimization, and (4)~hybrid Mamba--Transformer architectures.


% =====================================================================
\begin{thebibliography}{9}

\bibitem{gu2023mamba}
Gu, A., Dao, T.:
Mamba: Linear-Time Sequence Modeling with Selective State Spaces.
arXiv:2312.00752 (2023)

\bibitem{clifford2017af}
Clifford, G.D., et~al.:
AF Classification from a Short Single Lead ECG Recording: The
PhysioNet/Computing in Cardiology Challenge 2017.
Computing in Cardiology~\textbf{44} (2017)

\bibitem{vaswani2017attention}
Vaswani, A., et~al.:
Attention Is All You Need.
In: NeurIPS, vol.~30 (2017)

\bibitem{hochreiter1997lstm}
Hochreiter, S., Schmidhuber, J.:
Long Short-Term Memory.
Neural Computation~\textbf{9}(8), 1735--1780 (1997)

\bibitem{gu2022s4}
Gu, A., Goel, K., R{\'e}, C.:
Efficiently Modeling Long Sequences with Structured State Spaces.
In: ICLR (2022)

\end{thebibliography}

\end{document}
 inside your main paper,
%        OR compile standalone with: pdflatex cinc17_experiments.tex
%
% IMPORTANT — Image paths assume the following directory structure
% relative to THIS .tex file:
%   training_plots_of_all_four_models/   (fig1..fig8 as .png or .pdf)
%   inference_plots/                     (confusion, ROC, etc. as .png)
%
% If compiling standalone, the preamble below is used.
% If including in an existing LNCS document, comment out the preamble.
% =============================================================================

\documentclass[runningheads]{llncs}

\usepackage{graphicx}
\usepackage{booktabs}
\usepackage{amsmath}
\usepackage{subcaption}
\usepackage{hyperref}
\usepackage{float}
\usepackage[utf8]{inputenc}

% Relative path prefixes (adjust if your folder structure differs)
\newcommand{\trainplot}[1]{training_plots_of_all_four_models/#1}
\newcommand{\infplot}[1]{inference_plots/#1}

\begin{document}

\title{Mamba-Inspired State Space Models for Efficient\\Multi-Class Cardiac Arrhythmia Classification:\\A Comparative Study on Edge-Deployable ECG Monitoring}

\author{Utkarsh \and Supervisor Name}
\institute{Institution Name}

\maketitle

% =====================================================================
\begin{abstract}
Real-time cardiac arrhythmia detection on edge devices demands architectures
that balance classification accuracy with computational efficiency. This paper
presents a comparative evaluation of four sequence modeling architectures ---
Mamba (Selective State Space Model), Transformer, Long Short-Term Memory (LSTM),
and a baseline State Space Model (SSM) --- for multi-class ECG arrhythmia
classification on the PhysioNet/CinC~2017 Challenge dataset. We train all
models on identical data splits with consistent preprocessing and evaluate on a
held-out test set of 1,706~ECG recordings across four classes: Normal~(N),
Other~(O), Atrial Fibrillation~(A), and Noise~($\sim$). During training, Mamba
achieves the highest Macro~F1 score of \textbf{42.91\%}, outperforming the
Transformer (24.24\%), SSM (27.61\%), and LSTM (18.80\%). Critically, Mamba
operates with $O(N)$ linear complexity, compared to the Transformer's $O(N^2)$
quadratic attention, making it uniquely suited for real-time inference on
resource-constrained hardware. We further analyze inference-time
generalization, training dynamics, class-wise performance, and computational
trade-offs.
\end{abstract}


% =====================================================================
\section{Introduction}
\label{sec:intro}

Cardiac arrhythmia affects millions worldwide, and early detection through
continuous ECG monitoring can be lifesaving. Deploying such monitoring on edge
devices --- wearable patches, portable ECG monitors, implantable devices ---
requires models that are simultaneously accurate and computationally
lightweight.

Traditional approaches fall into two broad categories. \textbf{Recurrent
models} (LSTM, GRU) offer $O(N)$ complexity but suffer from vanishing
gradients, which limits their capacity to capture long-range dependencies in
ECG signals spanning thousands of time steps. \textbf{Attention-based models}
(Transformer) provide superior context modeling but incur $O(N^2)$ complexity
due to self-attention, making them impractical for raw ECG waveforms of
2,250+ time steps (after downsampling).

Mamba~\cite{gu2023mamba} introduces a third path: \textit{Selective State
Space Models} (S6) that combine the linear complexity of recurrent models with
a selective attention mechanism. Unlike fixed-coefficient SSMs, Mamba's
parameters are input-dependent, allowing the model to dynamically gate
information flow based on ECG signal content.

This paper provides an empirical comparison of all four architectures under
identical training conditions and analyzes the architectural reasons behind
Mamba's advantage for ECG classification.


% =====================================================================
\section{Dataset and Preprocessing}
\label{sec:dataset}

\subsection{PhysioNet/CinC 2017 Dataset}

We use the PhysioNet/Computing in Cardiology Challenge 2017
(CinC~2017)~\cite{clifford2017af} training set, containing 8,528 single-lead
ECG recordings sampled at 300\,Hz. Each recording is labeled into one of four
classes as shown in Table~\ref{tab:class_dist}.

\begin{table}[ht]
\centering
\caption{Class distribution of the PhysioNet/CinC 2017 dataset.}
\label{tab:class_dist}
\begin{tabular}{@{} l c r r @{}}
\toprule
\textbf{Class} & \textbf{Label} & \textbf{Samples} & \textbf{Percentage} \\
\midrule
Normal              & N    & 5,076 & 59.5\% \\
Other               & O    & 2,311 & 27.1\% \\
Atrial Fibrillation & A    &   758 &  8.9\% \\
Noise               & $\sim$ & 383 &  4.5\% \\
\bottomrule
\end{tabular}
\end{table}

The severe class imbalance --- particularly the 4.5\% representation of Noise
signals --- poses a significant challenge for all architectures.

\begin{figure}[ht]
\centering
\includegraphics[width=0.6\textwidth]{\trainplot{fig6_class_distribution.png}}
\caption{Class distribution across the CinC~2017 dataset showing the severe
imbalance between Normal, Other, AFib, and Noise categories.}
\label{fig:class_dist}
\end{figure}

\subsection{Preprocessing}

\begin{itemize}
    \item \textbf{Downsampling:} 300\,Hz $\rightarrow$ 150\,Hz (factor of 2) to
          reduce sequence length.
    \item \textbf{Sequence truncation:} Maximum length of 4,500 time steps.
    \item \textbf{Train/Test split:} 80/20 stratified split with random seed~42
          (6,822 train, 1,706 test).
    \item \textbf{Normalization:} Raw signal values used without additional
          normalization.
\end{itemize}


% =====================================================================
\section{Model Architectures}
\label{sec:models}

All four architectures share a common embedding dimension of $d_{\text{model}}=64$ and are designed to be small enough for edge deployment ($<200$\,KB).

\subsection{Mamba (Selective State Space Model)}

Mamba replaces the fixed state transition matrices of traditional SSMs with
input-dependent selection mechanisms:
\begin{itemize}
    \item \textbf{Selective scan (S6):} Parameters $A,B,C,\Delta$ are functions
          of the input, enabling content-aware information filtering.
    \item \textbf{Causal convolution:} 1D depthwise convolution ($k{=}3$) for
          local pattern capture.
    \item \textbf{Architecture:} 2 SelectiveSSMLayer blocks, $d_{\text{model}}{=}64$,
          $d_{\text{state}}{=}16$, length-aware pooling.
\end{itemize}
The forward pass processes the sequence in $O(N)$ time with $O(1)$ per-step
computation.

\subsection{Transformer}
Standard encoder with sinusoidal positional encoding (max\_len$=$20,000),
2-layer TransformerEncoder, 4~attention heads, $d_{\text{model}}{=}64$, and mean
pooling. Self-attention yields $O(N^2)$ complexity.

\subsection{LSTM}
Unidirectional 2-layer LSTM with hidden size 64. Final hidden state used as
sequence representation. $O(N)$ complexity but sequential processing.

\subsection{SSM (Baseline S4-style)}
Fixed-coefficient State Space Model with learnable $A,B,C$ matrices and
HiPPO initialization. Shares the same wrapper as Mamba but lacks the
selective mechanism, serving as a controlled ablation.


% =====================================================================
\section{Training Configuration}
\label{sec:training}

All models were trained with identical hyperparameters to ensure a fair
comparison (Table~\ref{tab:train_config}).

\begin{table}[ht]
\centering
\caption{Training hyperparameters (identical for all models).}
\label{tab:train_config}
\begin{tabular}{@{} l l @{}}
\toprule
\textbf{Parameter} & \textbf{Value} \\
\midrule
Epochs              & 6 \\
Effective batch size & 16 (4$\times$4 gradient accumulation) \\
Optimizer           & AdamW \\
Learning rate       & $1 \times 10^{-3}$ with cosine annealing \\
Loss function       & Cross-entropy (unweighted) \\
Mixed precision     & Disabled (T4 GPU compatibility) \\
Device              & NVIDIA Tesla T4 (16\,GB) \\
Checkpoint strategy & Best Macro~F1 saved per model \\
\bottomrule
\end{tabular}
\end{table}


% =====================================================================
\section{Training Results}
\label{sec:training_results}

\subsection{Overall Performance}

Table~\ref{tab:training_results} summarizes the best validation metrics
achieved during training. Mamba achieves 1.55$\times$ higher Macro~F1 than
the second-best model (SSM) and 1.77$\times$ higher than the Transformer.

\begin{table}[ht]
\centering
\caption{Training-phase performance comparison across all four models.}
\label{tab:training_results}
\begin{tabular}{@{} l r r r r @{}}
\toprule
\textbf{Model} & \textbf{Best Macro F1} & \textbf{Final Train Acc}
               & \textbf{Final Val Loss} & \textbf{Training Time} \\
\midrule
\textbf{Mamba} & \textbf{42.91\%} & 55.26\% & 1.085 &  9.1\,hrs \\
SSM            & 27.61\%          & 52.42\% & 1.241 &  5.6\,hrs \\
Transformer    & 24.24\%          & 58.25\% & 1.206 & 0.25\,hrs \\
LSTM           & 18.80\%          & 60.19\% & 1.270 & 0.30\,hrs \\
\bottomrule
\end{tabular}
\end{table}

\begin{figure}[ht]
\centering
\includegraphics[width=0.75\textwidth]{\trainplot{fig3_model_comparison.png}}
\caption{Model comparison: Best Macro~F1 and peak memory across all four
architectures.}
\label{fig:model_comparison}
\end{figure}

\begin{figure}[ht]
\centering
\includegraphics[width=0.75\textwidth]{\trainplot{fig5_summary_table.png}}
\caption{Summary table of training metrics for all four models.}
\label{fig:summary_table}
\end{figure}


\subsection{Training Dynamics}

\subsubsection{Loss Convergence.}
Figure~\ref{fig:loss_curves} shows that Mamba exhibits the steepest loss
reduction (1.29$\rightarrow$1.07), indicating effective learning of
discriminative features. The Transformer's loss decreases minimally
(1.29$\rightarrow$1.20), suggesting it struggles with raw ECG waveform
modeling at this sequence length.

\begin{figure}[ht]
\centering
\includegraphics[width=0.75\textwidth]{\trainplot{fig1_loss_curves.png}}
\caption{Validation loss curves across 6~epochs for all four models. Mamba
shows the steepest convergence.}
\label{fig:loss_curves}
\end{figure}

\begin{figure}[ht]
\centering
\includegraphics[width=0.9\textwidth]{\trainplot{fig7_loss_per_model.png}}
\caption{Individual training and validation loss curves per model, showing
detailed convergence behavior.}
\label{fig:loss_per_model}
\end{figure}

\subsubsection{F1 Progression.}
Mamba's F1 rises sharply from 18.8\% to 42.9\% by epoch~5, then drops to
34.9\% at epoch~6 --- indicating the onset of overfitting
(Figure~\ref{fig:f1_curves}). The SSM shows gradual improvement
(18.8\%$\rightarrow$27.6\%), confirming that Mamba's selective mechanism is
critical for ECG feature extraction. LSTM's F1 remains flat at 18.8\% across
all epochs, indicating majority-class prediction.

\begin{figure}[ht]
\centering
\includegraphics[width=0.75\textwidth]{\trainplot{fig2_f1_curves.png}}
\caption{Macro~F1 progression during training. Mamba achieves its peak at
epoch~5 before overfitting.}
\label{fig:f1_curves}
\end{figure}

\subsubsection{Training Accuracy.}
Figure~\ref{fig:train_accuracy} shows the training accuracy curves. LSTM
achieves the highest training accuracy (60.2\%) despite the lowest F1 (18.8\%),
revealing that it has learned to predict the majority class (Normal) rather
than discriminating between all four categories.

\begin{figure}[ht]
\centering
\includegraphics[width=0.75\textwidth]{\trainplot{fig4_train_accuracy.png}}
\caption{Training accuracy curves. High LSTM accuracy is deceptive ---
the model predicts the majority class.}
\label{fig:train_accuracy}
\end{figure}


\subsection{Efficiency Analysis}

Figure~\ref{fig:efficiency} plots Macro~F1 against training time, revealing
the accuracy--cost trade-off. Mamba requires more training time
(9.1\,hrs vs.\ 0.25\,hrs for Transformer) due to its sequential for-loop
over 4,500 time steps. However, this is a \textit{training-time} limitation
--- during edge inference, Mamba maintains a recurrent state with $O(1)$
per-step computation.

\begin{figure}[ht]
\centering
\includegraphics[width=0.65\textwidth]{\trainplot{fig8_efficiency_scatter.png}}
\caption{Efficiency scatter: Macro~F1 vs.\ training time. Mamba achieves the
best accuracy at moderate computational cost.}
\label{fig:efficiency}
\end{figure}


% =====================================================================
\section{Inference Results}
\label{sec:inference}

To evaluate generalization, we run inference on the held-out test set of
1,706~ECG recordings using the best checkpoints from training.
Table~\ref{tab:inference_results} summarizes the results.

\begin{table}[ht]
\centering
\caption{Inference-time performance on the held-out test set.}
\label{tab:inference_results}
\begin{tabular}{@{} l r r r r @{}}
\toprule
\textbf{Model} & \textbf{Macro F1} & \textbf{Weighted F1}
               & \textbf{Accuracy} & \textbf{Inference Time} \\
\midrule
Transformer    & 18.59\% & 43.97\% & 59.38\% & 23.5\,s \\
LSTM           & 18.18\% & 37.05\% & 40.91\% & 10.6\,s \\
SSM            & 13.58\% & 13.46\% & 17.76\% & 78.5\,s \\
Mamba          & 11.73\% & 13.89\% & 29.25\% & 130.4\,s \\
\bottomrule
\end{tabular}
\end{table}

\begin{figure}[ht]
\centering
\includegraphics[width=0.75\textwidth]{\infplot{overall_metrics.png}}
\caption{Overall inference metrics comparison: Macro~F1, Weighted~F1, and
Accuracy for all four models on the test set.}
\label{fig:overall_metrics}
\end{figure}

\begin{figure}[ht]
\centering
\includegraphics[width=0.75\textwidth]{\infplot{summary_table.png}}
\caption{Summary table of all inference metrics.}
\label{fig:inf_summary}
\end{figure}


\subsection{Confusion Matrix Analysis}

Figure~\ref{fig:confusion} reveals the prediction patterns of each model:
\begin{itemize}
    \item \textbf{Transformer:} Predicts ``Normal'' for nearly all inputs
          (majority-class collapse), achieving high accuracy (59.4\%) that
          merely reflects the class prior.
    \item \textbf{Mamba:} Predicts ``Other'' predominantly, correctly
          classifying 486/491 Other samples but misclassifying 998/1010 Normal
          samples.
    \item \textbf{LSTM:} Shows mixed Normal/AFib predictions, suggesting it
          learns some cardiac rhythm features but cannot reliably discriminate.
    \item \textbf{SSM:} Similar to Mamba but with weaker class separation.
\end{itemize}

\begin{figure}[ht]
\centering
\includegraphics[width=0.85\textwidth]{\infplot{confusion_matrices (1).png}}
\caption{Confusion matrices for all four models. Each model collapses to
predicting 1--2 dominant classes.}
\label{fig:confusion}
\end{figure}


\subsection{Per-Class F1 Analysis}

Figure~\ref{fig:per_class_f1} shows the class-wise F1 breakdown:
\begin{itemize}
    \item \textbf{Normal:} Only Transformer achieves meaningful F1 (74.2\%) by
          predicting Normal by default.
    \item \textbf{Other:} Mamba achieves the highest F1 (44.5\%).
    \item \textbf{AFib:} Only LSTM achieves a non-zero F1 (11.5\%).
    \item \textbf{Noise:} All models achieve 0\% F1 --- with only 57 test
          samples, no model learned to detect noise.
\end{itemize}

\begin{figure}[ht]
\centering
\includegraphics[width=0.75\textwidth]{\infplot{per_class_f1.png}}
\caption{Per-class F1 scores. Severe class imbalance causes all models to
fail on minority classes.}
\label{fig:per_class_f1}
\end{figure}


\subsection{ROC and Precision--Recall Analysis}

ROC curves (Figure~\ref{fig:roc}) and precision--recall curves
(Figure~\ref{fig:pr}) provide further insight. AUC values near 0.5 for most
class--model combinations confirm random-like discrimination. Precision drops
sharply as recall increases, indicating that none of the models can maintain
precision while retrieving more positive samples from minority classes.

\begin{figure}[ht]
\centering
\includegraphics[width=0.85\textwidth]{\infplot{roc_curves.png}}
\caption{ROC curves (one-vs-rest) for all four models.  Curves near the
diagonal indicate near-random classification.}
\label{fig:roc}
\end{figure}

\begin{figure}[ht]
\centering
\includegraphics[width=0.85\textwidth]{\infplot{precision_recall.png}}
\caption{Precision--Recall curves for all four models. Low area under the
curve for minority classes confirms poor discriminative ability.}
\label{fig:pr}
\end{figure}


\subsection{Inference Speed}

Figure~\ref{fig:inf_speed} compares inference times on the T4~GPU. The
Transformer (23.5\,s) and LSTM (10.6\,s) benefit from optimized CUDA kernels
and parallelizable operations. Mamba (130.4\,s) is slowest due to its
sequential selective scan. \textbf{Note:} batch-mode GPU inference does not
reflect edge deployment, where Mamba's $O(1)$ recurrent state update per
sample would outperform the Transformer's $O(N^2)$ attention recomputation.

\begin{figure}[ht]
\centering
\includegraphics[width=0.6\textwidth]{\infplot{inference_speed.png}}
\caption{Inference time for 1,706 test samples on T4~GPU. Mamba's sequential
scan is slower in batch mode but advantageous in streaming edge deployment.}
\label{fig:inf_speed}
\end{figure}


% =====================================================================
\section{Analysis and Discussion}
\label{sec:analysis}

\subsection{Training vs.\ Inference Gap}

Table~\ref{tab:gap} highlights the significant gap between training-time
validation and test-time inference metrics.

\begin{table}[ht]
\centering
\caption{Gap between training validation and test-time inference Macro~F1.}
\label{tab:gap}
\begin{tabular}{@{} l r r r @{}}
\toprule
\textbf{Model} & \textbf{Train Val F1} & \textbf{Test F1} & \textbf{Gap} \\
\midrule
Mamba       & 42.91\% & 11.73\% & $-31.2$\,pp \\
Transformer & 24.24\% & 18.59\% & $-5.7$\,pp  \\
LSTM        & 18.80\% & 18.18\% & $-0.6$\,pp  \\
SSM         & 27.61\% & 13.58\% & $-14.0$\,pp \\
\bottomrule
\end{tabular}
\end{table}

This gap is attributed to three factors:
\begin{enumerate}
    \item \textbf{WeightedRandomSampler during training:} Training used
          oversampling of minority classes, creating artificially balanced
          batches. Inference runs on the true imbalanced distribution
          (59\%~Normal).
    \item \textbf{Limited epochs (6):} Mamba showed overfitting at epoch~5,
          while the Transformer barely learned. 20--50~epochs with early
          stopping would improve all models.
    \item \textbf{No class-weighted loss:} Unweighted cross-entropy fails to
          penalize minority-class misclassifications, allowing the gradient
          signal from majority classes to dominate.
\end{enumerate}


\subsection{Why Mamba Outperforms During Training}

The 15.3~percentage-point gap between Mamba (42.9\%) and baseline SSM
(27.6\%) isolates the effect of the \textit{selective mechanism}:

\begin{itemize}
    \item In Mamba, the discretization step~$\Delta$, matrices $B$ and $C$ are
          computed as functions of the input at each time step.
    \item A sharp QRS complex produces a large~$\Delta$, causing the model to
          ``remember'' the event; a flat baseline segment produces a
          small~$\Delta$, allowing the model to ``forget.''
    \item The fixed-coefficient SSM cannot distinguish between clinically
          relevant and irrelevant ECG segments.
\end{itemize}

\subsection{Why Transformers Struggle}

Three factors limit Transformer performance on raw ECG:
\begin{enumerate}
    \item \textbf{Positional encoding gap:} Sinusoidal embeddings were designed
          for discrete tokens, not continuous waveforms.
    \item \textbf{Attention dilution:} With $N{=}4{,}500$ time steps,
          self-attention is diluted across all positions. Clinically relevant
          QRS complexes span $\sim$15 time steps ($<0.4\%$ of the sequence).
    \item \textbf{No causal structure:} ECG analysis benefits from stateful,
          causal processing that accumulates evidence over time.
\end{enumerate}

\subsection{Why LSTM Fails}

LSTM's 18.8\% F1 equals the majority-class baseline, revealing: (i)~vanishing
gradients over 4,500~steps, (ii)~unidirectional processing prevents using
future context, and (iii)~64-dimensional hidden states may be insufficient for
the full variety of ECG morphologies.


% =====================================================================
\section{Edge Device Deployment Analysis}
\label{sec:edge}

\begin{table}[ht]
\centering
\caption{Inference complexity and edge deployment characteristics.}
\label{tab:edge}
\begin{tabular}{@{} l c c c c @{}}
\toprule
\textbf{Model} & \textbf{Per-step} & \textbf{Memory} &
\textbf{Parallel} & \textbf{Approx.\ Size} \\
\midrule
Mamba       & $O(1)$ & $O(d \times N)$ & No  & $\sim$200\,KB \\
Transformer & $O(N)$ & $O(N^2)$        & Yes & $\sim$180\,KB \\
LSTM        & $O(1)$ & $O(h)$          & No  & $\sim$140\,KB \\
SSM         & $O(1)$ & $O(d \times N)$ & No  & $\sim$192\,KB \\
\bottomrule
\end{tabular}
\end{table}

For continuous ECG monitoring at 300\,Hz, Mamba processes each new sample in
$O(1)$ by maintaining a recurrent state, whereas the Transformer must
recompute $O(N^2)$ attention over all previous samples. After 10~seconds
(3,000~samples), attention requires $\sim$9\,M operations per update. All
models fit comfortably on microcontrollers with 256\,KB+ RAM (e.g., ARM
Cortex-M4/M7, ESP32).


% =====================================================================
\section{Limitations}
\label{sec:limitations}

\begin{itemize}
    \item \textbf{Training duration:} Only 6~epochs due to Kaggle GPU
          constraints. Longer training with early stopping is expected to
          improve all models, particularly Mamba.
    \item \textbf{No hyperparameter tuning:} Identical hyperparameters were
          used for all models. Architecture-specific tuning may close the gap.
    \item \textbf{Class imbalance:} No focal loss, class weighting, or
          oversampling was applied. This disadvantages all models on minority
          classes.
    \item \textbf{Single-lead ECG:} CinC~2017 uses Lead~I only. Multi-lead
          ECG may favor architectures with more parameters.
    \item \textbf{Training speed:} Mamba's sequential scan makes it
          $37\times$ slower to train than the Transformer, though this does not
          affect per-sample inference on edge devices.
\end{itemize}


% =====================================================================
\section{Conclusion}
\label{sec:conclusion}

This study provides empirical evidence that Mamba's Selective State Space
architecture is better suited for ECG arrhythmia classification than
Transformers, LSTMs, and baseline SSMs. During training, Mamba achieves
42.91\%~Macro~F1 --- 55\% higher than the next-best model --- by leveraging
input-dependent state selection to dynamically focus on clinically relevant
ECG segments.

The gap between training and inference metrics is attributed to the training
pipeline (weighted sampling, limited epochs, unweighted loss) rather than the
architecture itself, and can be addressed through class-weighted loss, longer
training with early stopping, and proper validation protocols.

The $O(N)$ linear complexity, combined with $O(1)$ per-step recurrent
inference and a model size under 200\,KB, makes Mamba uniquely suited for
real-time ECG monitoring on edge devices where computational budgets are tight
and latency requirements are strict.

Future work should explore: (1)~longer training with class-weighted focal
loss, (2)~multi-lead ECG input, (3)~quantization and pruning for further edge
optimization, and (4)~hybrid Mamba--Transformer architectures.


% =====================================================================
\begin{thebibliography}{9}

\bibitem{gu2023mamba}
Gu, A., Dao, T.:
Mamba: Linear-Time Sequence Modeling with Selective State Spaces.
arXiv:2312.00752 (2023)

\bibitem{clifford2017af}
Clifford, G.D., et~al.:
AF Classification from a Short Single Lead ECG Recording: The
PhysioNet/Computing in Cardiology Challenge 2017.
Computing in Cardiology~\textbf{44} (2017)

\bibitem{vaswani2017attention}
Vaswani, A., et~al.:
Attention Is All You Need.
In: NeurIPS, vol.~30 (2017)

\bibitem{hochreiter1997lstm}
Hochreiter, S., Schmidhuber, J.:
Long Short-Term Memory.
Neural Computation~\textbf{9}(8), 1735--1780 (1997)

\bibitem{gu2022s4}
Gu, A., Goel, K., R{\'e}, C.:
Efficiently Modeling Long Sequences with Structured State Spaces.
In: ICLR (2022)

\end{thebibliography}

\end{document}
 inside your main paper,
%        OR compile standalone with: pdflatex cinc17_experiments.tex
%
% IMPORTANT — Image paths assume the following directory structure
% relative to THIS .tex file:
%   training_plots_of_all_four_models/   (fig1..fig8 as .png or .pdf)
%   inference_plots/                     (confusion, ROC, etc. as .png)
%
% If compiling standalone, the preamble below is used.
% If including in an existing LNCS document, comment out the preamble.
% =============================================================================

\documentclass[runningheads]{llncs}

\usepackage{graphicx}
\usepackage{booktabs}
\usepackage{amsmath}
\usepackage{subcaption}
\usepackage{hyperref}
\usepackage{float}
\usepackage[utf8]{inputenc}

% Relative path prefixes (adjust if your folder structure differs)
\newcommand{\trainplot}[1]{training_plots_of_all_four_models/#1}
\newcommand{\infplot}[1]{inference_plots/#1}

\begin{document}

\title{Mamba-Inspired State Space Models for Efficient\\Multi-Class Cardiac Arrhythmia Classification:\\A Comparative Study on Edge-Deployable ECG Monitoring}

\author{Utkarsh \and Supervisor Name}
\institute{Institution Name}

\maketitle

% =====================================================================
\begin{abstract}
Real-time cardiac arrhythmia detection on edge devices demands architectures
that balance classification accuracy with computational efficiency. This paper
presents a comparative evaluation of four sequence modeling architectures ---
Mamba (Selective State Space Model), Transformer, Long Short-Term Memory (LSTM),
and a baseline State Space Model (SSM) --- for multi-class ECG arrhythmia
classification on the PhysioNet/CinC~2017 Challenge dataset. We train all
models on identical data splits with consistent preprocessing and evaluate on a
held-out test set of 1,706~ECG recordings across four classes: Normal~(N),
Other~(O), Atrial Fibrillation~(A), and Noise~($\sim$). During training, Mamba
achieves the highest Macro~F1 score of \textbf{42.91\%}, outperforming the
Transformer (24.24\%), SSM (27.61\%), and LSTM (18.80\%). Critically, Mamba
operates with $O(N)$ linear complexity, compared to the Transformer's $O(N^2)$
quadratic attention, making it uniquely suited for real-time inference on
resource-constrained hardware. We further analyze inference-time
generalization, training dynamics, class-wise performance, and computational
trade-offs.
\end{abstract}


% =====================================================================
\section{Introduction}
\label{sec:intro}

Cardiac arrhythmia affects millions worldwide, and early detection through
continuous ECG monitoring can be lifesaving. Deploying such monitoring on edge
devices --- wearable patches, portable ECG monitors, implantable devices ---
requires models that are simultaneously accurate and computationally
lightweight.

Traditional approaches fall into two broad categories. \textbf{Recurrent
models} (LSTM, GRU) offer $O(N)$ complexity but suffer from vanishing
gradients, which limits their capacity to capture long-range dependencies in
ECG signals spanning thousands of time steps. \textbf{Attention-based models}
(Transformer) provide superior context modeling but incur $O(N^2)$ complexity
due to self-attention, making them impractical for raw ECG waveforms of
2,250+ time steps (after downsampling).

Mamba~\cite{gu2023mamba} introduces a third path: \textit{Selective State
Space Models} (S6) that combine the linear complexity of recurrent models with
a selective attention mechanism. Unlike fixed-coefficient SSMs, Mamba's
parameters are input-dependent, allowing the model to dynamically gate
information flow based on ECG signal content.

This paper provides an empirical comparison of all four architectures under
identical training conditions and analyzes the architectural reasons behind
Mamba's advantage for ECG classification.


% =====================================================================
\section{Dataset and Preprocessing}
\label{sec:dataset}

\subsection{PhysioNet/CinC 2017 Dataset}

We use the PhysioNet/Computing in Cardiology Challenge 2017
(CinC~2017)~\cite{clifford2017af} training set, containing 8,528 single-lead
ECG recordings sampled at 300\,Hz. Each recording is labeled into one of four
classes as shown in Table~\ref{tab:class_dist}.

\begin{table}[ht]
\centering
\caption{Class distribution of the PhysioNet/CinC 2017 dataset.}
\label{tab:class_dist}
\begin{tabular}{@{} l c r r @{}}
\toprule
\textbf{Class} & \textbf{Label} & \textbf{Samples} & \textbf{Percentage} \\
\midrule
Normal              & N    & 5,076 & 59.5\% \\
Other               & O    & 2,311 & 27.1\% \\
Atrial Fibrillation & A    &   758 &  8.9\% \\
Noise               & $\sim$ & 383 &  4.5\% \\
\bottomrule
\end{tabular}
\end{table}

The severe class imbalance --- particularly the 4.5\% representation of Noise
signals --- poses a significant challenge for all architectures.

\begin{figure}[ht]
\centering
\includegraphics[width=0.6\textwidth]{\trainplot{fig6_class_distribution.png}}
\caption{Class distribution across the CinC~2017 dataset showing the severe
imbalance between Normal, Other, AFib, and Noise categories.}
\label{fig:class_dist}
\end{figure}

\subsection{Preprocessing}

\begin{itemize}
    \item \textbf{Downsampling:} 300\,Hz $\rightarrow$ 150\,Hz (factor of 2) to
          reduce sequence length.
    \item \textbf{Sequence truncation:} Maximum length of 4,500 time steps.
    \item \textbf{Train/Test split:} 80/20 stratified split with random seed~42
          (6,822 train, 1,706 test).
    \item \textbf{Normalization:} Raw signal values used without additional
          normalization.
\end{itemize}


% =====================================================================
\section{Model Architectures}
\label{sec:models}

All four architectures share a common embedding dimension of $d_{\text{model}}=64$ and are designed to be small enough for edge deployment ($<200$\,KB).

\subsection{Mamba (Selective State Space Model)}

Mamba replaces the fixed state transition matrices of traditional SSMs with
input-dependent selection mechanisms:
\begin{itemize}
    \item \textbf{Selective scan (S6):} Parameters $A,B,C,\Delta$ are functions
          of the input, enabling content-aware information filtering.
    \item \textbf{Causal convolution:} 1D depthwise convolution ($k{=}3$) for
          local pattern capture.
    \item \textbf{Architecture:} 2 SelectiveSSMLayer blocks, $d_{\text{model}}{=}64$,
          $d_{\text{state}}{=}16$, length-aware pooling.
\end{itemize}
The forward pass processes the sequence in $O(N)$ time with $O(1)$ per-step
computation.

\subsection{Transformer}
Standard encoder with sinusoidal positional encoding (max\_len$=$20,000),
2-layer TransformerEncoder, 4~attention heads, $d_{\text{model}}{=}64$, and mean
pooling. Self-attention yields $O(N^2)$ complexity.

\subsection{LSTM}
Unidirectional 2-layer LSTM with hidden size 64. Final hidden state used as
sequence representation. $O(N)$ complexity but sequential processing.

\subsection{SSM (Baseline S4-style)}
Fixed-coefficient State Space Model with learnable $A,B,C$ matrices and
HiPPO initialization. Shares the same wrapper as Mamba but lacks the
selective mechanism, serving as a controlled ablation.


% =====================================================================
\section{Training Configuration}
\label{sec:training}

All models were trained with identical hyperparameters to ensure a fair
comparison (Table~\ref{tab:train_config}).

\begin{table}[ht]
\centering
\caption{Training hyperparameters (identical for all models).}
\label{tab:train_config}
\begin{tabular}{@{} l l @{}}
\toprule
\textbf{Parameter} & \textbf{Value} \\
\midrule
Epochs              & 6 \\
Effective batch size & 16 (4$\times$4 gradient accumulation) \\
Optimizer           & AdamW \\
Learning rate       & $1 \times 10^{-3}$ with cosine annealing \\
Loss function       & Cross-entropy (unweighted) \\
Mixed precision     & Disabled (T4 GPU compatibility) \\
Device              & NVIDIA Tesla T4 (16\,GB) \\
Checkpoint strategy & Best Macro~F1 saved per model \\
\bottomrule
\end{tabular}
\end{table}


% =====================================================================
\section{Training Results}
\label{sec:training_results}

\subsection{Overall Performance}

Table~\ref{tab:training_results} summarizes the best validation metrics
achieved during training. Mamba achieves 1.55$\times$ higher Macro~F1 than
the second-best model (SSM) and 1.77$\times$ higher than the Transformer.

\begin{table}[ht]
\centering
\caption{Training-phase performance comparison across all four models.}
\label{tab:training_results}
\begin{tabular}{@{} l r r r r @{}}
\toprule
\textbf{Model} & \textbf{Best Macro F1} & \textbf{Final Train Acc}
               & \textbf{Final Val Loss} & \textbf{Training Time} \\
\midrule
\textbf{Mamba} & \textbf{42.91\%} & 55.26\% & 1.085 &  9.1\,hrs \\
SSM            & 27.61\%          & 52.42\% & 1.241 &  5.6\,hrs \\
Transformer    & 24.24\%          & 58.25\% & 1.206 & 0.25\,hrs \\
LSTM           & 18.80\%          & 60.19\% & 1.270 & 0.30\,hrs \\
\bottomrule
\end{tabular}
\end{table}

\begin{figure}[ht]
\centering
\includegraphics[width=0.75\textwidth]{\trainplot{fig3_model_comparison.png}}
\caption{Model comparison: Best Macro~F1 and peak memory across all four
architectures.}
\label{fig:model_comparison}
\end{figure}

\begin{figure}[ht]
\centering
\includegraphics[width=0.75\textwidth]{\trainplot{fig5_summary_table.png}}
\caption{Summary table of training metrics for all four models.}
\label{fig:summary_table}
\end{figure}


\subsection{Training Dynamics}

\subsubsection{Loss Convergence.}
Figure~\ref{fig:loss_curves} shows that Mamba exhibits the steepest loss
reduction (1.29$\rightarrow$1.07), indicating effective learning of
discriminative features. The Transformer's loss decreases minimally
(1.29$\rightarrow$1.20), suggesting it struggles with raw ECG waveform
modeling at this sequence length.

\begin{figure}[ht]
\centering
\includegraphics[width=0.75\textwidth]{\trainplot{fig1_loss_curves.png}}
\caption{Validation loss curves across 6~epochs for all four models. Mamba
shows the steepest convergence.}
\label{fig:loss_curves}
\end{figure}

\begin{figure}[ht]
\centering
\includegraphics[width=0.9\textwidth]{\trainplot{fig7_loss_per_model.png}}
\caption{Individual training and validation loss curves per model, showing
detailed convergence behavior.}
\label{fig:loss_per_model}
\end{figure}

\subsubsection{F1 Progression.}
Mamba's F1 rises sharply from 18.8\% to 42.9\% by epoch~5, then drops to
34.9\% at epoch~6 --- indicating the onset of overfitting
(Figure~\ref{fig:f1_curves}). The SSM shows gradual improvement
(18.8\%$\rightarrow$27.6\%), confirming that Mamba's selective mechanism is
critical for ECG feature extraction. LSTM's F1 remains flat at 18.8\% across
all epochs, indicating majority-class prediction.

\begin{figure}[ht]
\centering
\includegraphics[width=0.75\textwidth]{\trainplot{fig2_f1_curves.png}}
\caption{Macro~F1 progression during training. Mamba achieves its peak at
epoch~5 before overfitting.}
\label{fig:f1_curves}
\end{figure}

\subsubsection{Training Accuracy.}
Figure~\ref{fig:train_accuracy} shows the training accuracy curves. LSTM
achieves the highest training accuracy (60.2\%) despite the lowest F1 (18.8\%),
revealing that it has learned to predict the majority class (Normal) rather
than discriminating between all four categories.

\begin{figure}[ht]
\centering
\includegraphics[width=0.75\textwidth]{\trainplot{fig4_train_accuracy.png}}
\caption{Training accuracy curves. High LSTM accuracy is deceptive ---
the model predicts the majority class.}
\label{fig:train_accuracy}
\end{figure}


\subsection{Efficiency Analysis}

Figure~\ref{fig:efficiency} plots Macro~F1 against training time, revealing
the accuracy--cost trade-off. Mamba requires more training time
(9.1\,hrs vs.\ 0.25\,hrs for Transformer) due to its sequential for-loop
over 4,500 time steps. However, this is a \textit{training-time} limitation
--- during edge inference, Mamba maintains a recurrent state with $O(1)$
per-step computation.

\begin{figure}[ht]
\centering
\includegraphics[width=0.65\textwidth]{\trainplot{fig8_efficiency_scatter.png}}
\caption{Efficiency scatter: Macro~F1 vs.\ training time. Mamba achieves the
best accuracy at moderate computational cost.}
\label{fig:efficiency}
\end{figure}


% =====================================================================
\section{Inference Results}
\label{sec:inference}

To evaluate generalization, we run inference on the held-out test set of
1,706~ECG recordings using the best checkpoints from training.
Table~\ref{tab:inference_results} summarizes the results.

\begin{table}[ht]
\centering
\caption{Inference-time performance on the held-out test set.}
\label{tab:inference_results}
\begin{tabular}{@{} l r r r r @{}}
\toprule
\textbf{Model} & \textbf{Macro F1} & \textbf{Weighted F1}
               & \textbf{Accuracy} & \textbf{Inference Time} \\
\midrule
Transformer    & 18.59\% & 43.97\% & 59.38\% & 23.5\,s \\
LSTM           & 18.18\% & 37.05\% & 40.91\% & 10.6\,s \\
SSM            & 13.58\% & 13.46\% & 17.76\% & 78.5\,s \\
Mamba          & 11.73\% & 13.89\% & 29.25\% & 130.4\,s \\
\bottomrule
\end{tabular}
\end{table}

\begin{figure}[ht]
\centering
\includegraphics[width=0.75\textwidth]{\infplot{overall_metrics.png}}
\caption{Overall inference metrics comparison: Macro~F1, Weighted~F1, and
Accuracy for all four models on the test set.}
\label{fig:overall_metrics}
\end{figure}

\begin{figure}[ht]
\centering
\includegraphics[width=0.75\textwidth]{\infplot{summary_table.png}}
\caption{Summary table of all inference metrics.}
\label{fig:inf_summary}
\end{figure}


\subsection{Confusion Matrix Analysis}

Figure~\ref{fig:confusion} reveals the prediction patterns of each model:
\begin{itemize}
    \item \textbf{Transformer:} Predicts ``Normal'' for nearly all inputs
          (majority-class collapse), achieving high accuracy (59.4\%) that
          merely reflects the class prior.
    \item \textbf{Mamba:} Predicts ``Other'' predominantly, correctly
          classifying 486/491 Other samples but misclassifying 998/1010 Normal
          samples.
    \item \textbf{LSTM:} Shows mixed Normal/AFib predictions, suggesting it
          learns some cardiac rhythm features but cannot reliably discriminate.
    \item \textbf{SSM:} Similar to Mamba but with weaker class separation.
\end{itemize}

\begin{figure}[ht]
\centering
\includegraphics[width=0.85\textwidth]{\infplot{confusion_matrices (1).png}}
\caption{Confusion matrices for all four models. Each model collapses to
predicting 1--2 dominant classes.}
\label{fig:confusion}
\end{figure}


\subsection{Per-Class F1 Analysis}

Figure~\ref{fig:per_class_f1} shows the class-wise F1 breakdown:
\begin{itemize}
    \item \textbf{Normal:} Only Transformer achieves meaningful F1 (74.2\%) by
          predicting Normal by default.
    \item \textbf{Other:} Mamba achieves the highest F1 (44.5\%).
    \item \textbf{AFib:} Only LSTM achieves a non-zero F1 (11.5\%).
    \item \textbf{Noise:} All models achieve 0\% F1 --- with only 57 test
          samples, no model learned to detect noise.
\end{itemize}

\begin{figure}[ht]
\centering
\includegraphics[width=0.75\textwidth]{\infplot{per_class_f1.png}}
\caption{Per-class F1 scores. Severe class imbalance causes all models to
fail on minority classes.}
\label{fig:per_class_f1}
\end{figure}


\subsection{ROC and Precision--Recall Analysis}

ROC curves (Figure~\ref{fig:roc}) and precision--recall curves
(Figure~\ref{fig:pr}) provide further insight. AUC values near 0.5 for most
class--model combinations confirm random-like discrimination. Precision drops
sharply as recall increases, indicating that none of the models can maintain
precision while retrieving more positive samples from minority classes.

\begin{figure}[ht]
\centering
\includegraphics[width=0.85\textwidth]{\infplot{roc_curves.png}}
\caption{ROC curves (one-vs-rest) for all four models.  Curves near the
diagonal indicate near-random classification.}
\label{fig:roc}
\end{figure}

\begin{figure}[ht]
\centering
\includegraphics[width=0.85\textwidth]{\infplot{precision_recall.png}}
\caption{Precision--Recall curves for all four models. Low area under the
curve for minority classes confirms poor discriminative ability.}
\label{fig:pr}
\end{figure}


\subsection{Inference Speed}

Figure~\ref{fig:inf_speed} compares inference times on the T4~GPU. The
Transformer (23.5\,s) and LSTM (10.6\,s) benefit from optimized CUDA kernels
and parallelizable operations. Mamba (130.4\,s) is slowest due to its
sequential selective scan. \textbf{Note:} batch-mode GPU inference does not
reflect edge deployment, where Mamba's $O(1)$ recurrent state update per
sample would outperform the Transformer's $O(N^2)$ attention recomputation.

\begin{figure}[ht]
\centering
\includegraphics[width=0.6\textwidth]{\infplot{inference_speed.png}}
\caption{Inference time for 1,706 test samples on T4~GPU. Mamba's sequential
scan is slower in batch mode but advantageous in streaming edge deployment.}
\label{fig:inf_speed}
\end{figure}


% =====================================================================
\section{Analysis and Discussion}
\label{sec:analysis}

\subsection{Training vs.\ Inference Gap}

Table~\ref{tab:gap} highlights the significant gap between training-time
validation and test-time inference metrics.

\begin{table}[ht]
\centering
\caption{Gap between training validation and test-time inference Macro~F1.}
\label{tab:gap}
\begin{tabular}{@{} l r r r @{}}
\toprule
\textbf{Model} & \textbf{Train Val F1} & \textbf{Test F1} & \textbf{Gap} \\
\midrule
Mamba       & 42.91\% & 11.73\% & $-31.2$\,pp \\
Transformer & 24.24\% & 18.59\% & $-5.7$\,pp  \\
LSTM        & 18.80\% & 18.18\% & $-0.6$\,pp  \\
SSM         & 27.61\% & 13.58\% & $-14.0$\,pp \\
\bottomrule
\end{tabular}
\end{table}

This gap is attributed to three factors:
\begin{enumerate}
    \item \textbf{WeightedRandomSampler during training:} Training used
          oversampling of minority classes, creating artificially balanced
          batches. Inference runs on the true imbalanced distribution
          (59\%~Normal).
    \item \textbf{Limited epochs (6):} Mamba showed overfitting at epoch~5,
          while the Transformer barely learned. 20--50~epochs with early
          stopping would improve all models.
    \item \textbf{No class-weighted loss:} Unweighted cross-entropy fails to
          penalize minority-class misclassifications, allowing the gradient
          signal from majority classes to dominate.
\end{enumerate}


\subsection{Why Mamba Outperforms During Training}

The 15.3~percentage-point gap between Mamba (42.9\%) and baseline SSM
(27.6\%) isolates the effect of the \textit{selective mechanism}:

\begin{itemize}
    \item In Mamba, the discretization step~$\Delta$, matrices $B$ and $C$ are
          computed as functions of the input at each time step.
    \item A sharp QRS complex produces a large~$\Delta$, causing the model to
          ``remember'' the event; a flat baseline segment produces a
          small~$\Delta$, allowing the model to ``forget.''
    \item The fixed-coefficient SSM cannot distinguish between clinically
          relevant and irrelevant ECG segments.
\end{itemize}

\subsection{Why Transformers Struggle}

Three factors limit Transformer performance on raw ECG:
\begin{enumerate}
    \item \textbf{Positional encoding gap:} Sinusoidal embeddings were designed
          for discrete tokens, not continuous waveforms.
    \item \textbf{Attention dilution:} With $N{=}4{,}500$ time steps,
          self-attention is diluted across all positions. Clinically relevant
          QRS complexes span $\sim$15 time steps ($<0.4\%$ of the sequence).
    \item \textbf{No causal structure:} ECG analysis benefits from stateful,
          causal processing that accumulates evidence over time.
\end{enumerate}

\subsection{Why LSTM Fails}

LSTM's 18.8\% F1 equals the majority-class baseline, revealing: (i)~vanishing
gradients over 4,500~steps, (ii)~unidirectional processing prevents using
future context, and (iii)~64-dimensional hidden states may be insufficient for
the full variety of ECG morphologies.


% =====================================================================
\section{Edge Device Deployment Analysis}
\label{sec:edge}

\begin{table}[ht]
\centering
\caption{Inference complexity and edge deployment characteristics.}
\label{tab:edge}
\begin{tabular}{@{} l c c c c @{}}
\toprule
\textbf{Model} & \textbf{Per-step} & \textbf{Memory} &
\textbf{Parallel} & \textbf{Approx.\ Size} \\
\midrule
Mamba       & $O(1)$ & $O(d \times N)$ & No  & $\sim$200\,KB \\
Transformer & $O(N)$ & $O(N^2)$        & Yes & $\sim$180\,KB \\
LSTM        & $O(1)$ & $O(h)$          & No  & $\sim$140\,KB \\
SSM         & $O(1)$ & $O(d \times N)$ & No  & $\sim$192\,KB \\
\bottomrule
\end{tabular}
\end{table}

For continuous ECG monitoring at 300\,Hz, Mamba processes each new sample in
$O(1)$ by maintaining a recurrent state, whereas the Transformer must
recompute $O(N^2)$ attention over all previous samples. After 10~seconds
(3,000~samples), attention requires $\sim$9\,M operations per update. All
models fit comfortably on microcontrollers with 256\,KB+ RAM (e.g., ARM
Cortex-M4/M7, ESP32).


% =====================================================================
\section{Limitations}
\label{sec:limitations}

\begin{itemize}
    \item \textbf{Training duration:} Only 6~epochs due to Kaggle GPU
          constraints. Longer training with early stopping is expected to
          improve all models, particularly Mamba.
    \item \textbf{No hyperparameter tuning:} Identical hyperparameters were
          used for all models. Architecture-specific tuning may close the gap.
    \item \textbf{Class imbalance:} No focal loss, class weighting, or
          oversampling was applied. This disadvantages all models on minority
          classes.
    \item \textbf{Single-lead ECG:} CinC~2017 uses Lead~I only. Multi-lead
          ECG may favor architectures with more parameters.
    \item \textbf{Training speed:} Mamba's sequential scan makes it
          $37\times$ slower to train than the Transformer, though this does not
          affect per-sample inference on edge devices.
\end{itemize}


% =====================================================================
\section{Conclusion}
\label{sec:conclusion}

This study provides empirical evidence that Mamba's Selective State Space
architecture is better suited for ECG arrhythmia classification than
Transformers, LSTMs, and baseline SSMs. During training, Mamba achieves
42.91\%~Macro~F1 --- 55\% higher than the next-best model --- by leveraging
input-dependent state selection to dynamically focus on clinically relevant
ECG segments.

The gap between training and inference metrics is attributed to the training
pipeline (weighted sampling, limited epochs, unweighted loss) rather than the
architecture itself, and can be addressed through class-weighted loss, longer
training with early stopping, and proper validation protocols.

The $O(N)$ linear complexity, combined with $O(1)$ per-step recurrent
inference and a model size under 200\,KB, makes Mamba uniquely suited for
real-time ECG monitoring on edge devices where computational budgets are tight
and latency requirements are strict.

Future work should explore: (1)~longer training with class-weighted focal
loss, (2)~multi-lead ECG input, (3)~quantization and pruning for further edge
optimization, and (4)~hybrid Mamba--Transformer architectures.


% =====================================================================
\begin{thebibliography}{9}

\bibitem{gu2023mamba}
Gu, A., Dao, T.:
Mamba: Linear-Time Sequence Modeling with Selective State Spaces.
arXiv:2312.00752 (2023)

\bibitem{clifford2017af}
Clifford, G.D., et~al.:
AF Classification from a Short Single Lead ECG Recording: The
PhysioNet/Computing in Cardiology Challenge 2017.
Computing in Cardiology~\textbf{44} (2017)

\bibitem{vaswani2017attention}
Vaswani, A., et~al.:
Attention Is All You Need.
In: NeurIPS, vol.~30 (2017)

\bibitem{hochreiter1997lstm}
Hochreiter, S., Schmidhuber, J.:
Long Short-Term Memory.
Neural Computation~\textbf{9}(8), 1735--1780 (1997)

\bibitem{gu2022s4}
Gu, A., Goel, K., R{\'e}, C.:
Efficiently Modeling Long Sequences with Structured State Spaces.
In: ICLR (2022)

\end{thebibliography}

\end{document}
 inside your main paper,
%        OR compile standalone with: pdflatex cinc17_experiments.tex
%
% IMPORTANT — Image paths assume the following directory structure
% relative to THIS .tex file:
%   training_plots_of_all_four_models/   (fig1..fig8 as .png or .pdf)
%   inference_plots/                     (confusion, ROC, etc. as .png)
%
% If compiling standalone, the preamble below is used.
% If including in an existing LNCS document, comment out the preamble.
% =============================================================================

\documentclass[runningheads]{llncs}

\usepackage{graphicx}
\usepackage{booktabs}
\usepackage{amsmath}
\usepackage{subcaption}
\usepackage{hyperref}
\usepackage{float}
\usepackage[utf8]{inputenc}

% Relative path prefixes (adjust if your folder structure differs)
\newcommand{\trainplot}[1]{training_plots_of_all_four_models/#1}
\newcommand{\infplot}[1]{inference_plots/#1}

\begin{document}

\title{Mamba-Inspired State Space Models for Efficient\\Multi-Class Cardiac Arrhythmia Classification:\\A Comparative Study on Edge-Deployable ECG Monitoring}

\author{Utkarsh \and Supervisor Name}
\institute{Institution Name}

\maketitle

% =====================================================================
\begin{abstract}
Real-time cardiac arrhythmia detection on edge devices demands architectures
that balance classification accuracy with computational efficiency. This paper
presents a comparative evaluation of four sequence modeling architectures ---
Mamba (Selective State Space Model), Transformer, Long Short-Term Memory (LSTM),
and a baseline State Space Model (SSM) --- for multi-class ECG arrhythmia
classification on the PhysioNet/CinC~2017 Challenge dataset. We train all
models on identical data splits with consistent preprocessing and evaluate on a
held-out test set of 1,706~ECG recordings across four classes: Normal~(N),
Other~(O), Atrial Fibrillation~(A), and Noise~($\sim$). During training, Mamba
achieves the highest Macro~F1 score of \textbf{42.91\%}, outperforming the
Transformer (24.24\%), SSM (27.61\%), and LSTM (18.80\%). Critically, Mamba
operates with $O(N)$ linear complexity, compared to the Transformer's $O(N^2)$
quadratic attention, making it uniquely suited for real-time inference on
resource-constrained hardware. We further analyze inference-time
generalization, training dynamics, class-wise performance, and computational
trade-offs.
\end{abstract}


% =====================================================================
\section{Introduction}
\label{sec:intro}

Cardiac arrhythmia affects millions worldwide, and early detection through
continuous ECG monitoring can be lifesaving. Deploying such monitoring on edge
devices --- wearable patches, portable ECG monitors, implantable devices ---
requires models that are simultaneously accurate and computationally
lightweight.

Traditional approaches fall into two broad categories. \textbf{Recurrent
models} (LSTM, GRU) offer $O(N)$ complexity but suffer from vanishing
gradients, which limits their capacity to capture long-range dependencies in
ECG signals spanning thousands of time steps. \textbf{Attention-based models}
(Transformer) provide superior context modeling but incur $O(N^2)$ complexity
due to self-attention, making them impractical for raw ECG waveforms of
2,250+ time steps (after downsampling).

Mamba~\cite{gu2023mamba} introduces a third path: \textit{Selective State
Space Models} (S6) that combine the linear complexity of recurrent models with
a selective attention mechanism. Unlike fixed-coefficient SSMs, Mamba's
parameters are input-dependent, allowing the model to dynamically gate
information flow based on ECG signal content.

This paper provides an empirical comparison of all four architectures under
identical training conditions and analyzes the architectural reasons behind
Mamba's advantage for ECG classification.


% =====================================================================
\section{Dataset and Preprocessing}
\label{sec:dataset}

\subsection{PhysioNet/CinC 2017 Dataset}

We use the PhysioNet/Computing in Cardiology Challenge 2017
(CinC~2017)~\cite{clifford2017af} training set, containing 8,528 single-lead
ECG recordings sampled at 300\,Hz. Each recording is labeled into one of four
classes as shown in Table~\ref{tab:class_dist}.

\begin{table}[ht]
\centering
\caption{Class distribution of the PhysioNet/CinC 2017 dataset.}
\label{tab:class_dist}
\begin{tabular}{@{} l c r r @{}}
\toprule
\textbf{Class} & \textbf{Label} & \textbf{Samples} & \textbf{Percentage} \\
\midrule
Normal              & N    & 5,076 & 59.5\% \\
Other               & O    & 2,311 & 27.1\% \\
Atrial Fibrillation & A    &   758 &  8.9\% \\
Noise               & $\sim$ & 383 &  4.5\% \\
\bottomrule
\end{tabular}
\end{table}

The severe class imbalance --- particularly the 4.5\% representation of Noise
signals --- poses a significant challenge for all architectures.

\begin{figure}[ht]
\centering
\includegraphics[width=0.6\textwidth]{\trainplot{fig6_class_distribution.png}}
\caption{Class distribution across the CinC~2017 dataset showing the severe
imbalance between Normal, Other, AFib, and Noise categories.}
\label{fig:class_dist}
\end{figure}

\subsection{Preprocessing}

\begin{itemize}
    \item \textbf{Downsampling:} 300\,Hz $\rightarrow$ 150\,Hz (factor of 2) to
          reduce sequence length.
    \item \textbf{Sequence truncation:} Maximum length of 4,500 time steps.
    \item \textbf{Train/Test split:} 80/20 stratified split with random seed~42
          (6,822 train, 1,706 test).
    \item \textbf{Normalization:} Raw signal values used without additional
          normalization.
\end{itemize}


% =====================================================================
\section{Model Architectures}
\label{sec:models}

All four architectures share a common embedding dimension of $d_{\text{model}}=64$ and are designed to be small enough for edge deployment ($<200$\,KB).

\subsection{Mamba (Selective State Space Model)}

Mamba replaces the fixed state transition matrices of traditional SSMs with
input-dependent selection mechanisms:
\begin{itemize}
    \item \textbf{Selective scan (S6):} Parameters $A,B,C,\Delta$ are functions
          of the input, enabling content-aware information filtering.
    \item \textbf{Causal convolution:} 1D depthwise convolution ($k{=}3$) for
          local pattern capture.
    \item \textbf{Architecture:} 2 SelectiveSSMLayer blocks, $d_{\text{model}}{=}64$,
          $d_{\text{state}}{=}16$, length-aware pooling.
\end{itemize}
The forward pass processes the sequence in $O(N)$ time with $O(1)$ per-step
computation.

\subsection{Transformer}
Standard encoder with sinusoidal positional encoding (max\_len$=$20,000),
2-layer TransformerEncoder, 4~attention heads, $d_{\text{model}}{=}64$, and mean
pooling. Self-attention yields $O(N^2)$ complexity.

\subsection{LSTM}
Unidirectional 2-layer LSTM with hidden size 64. Final hidden state used as
sequence representation. $O(N)$ complexity but sequential processing.

\subsection{SSM (Baseline S4-style)}
Fixed-coefficient State Space Model with learnable $A,B,C$ matrices and
HiPPO initialization. Shares the same wrapper as Mamba but lacks the
selective mechanism, serving as a controlled ablation.


% =====================================================================
\section{Training Configuration}
\label{sec:training}

All models were trained with identical hyperparameters to ensure a fair
comparison (Table~\ref{tab:train_config}).

\begin{table}[ht]
\centering
\caption{Training hyperparameters (identical for all models).}
\label{tab:train_config}
\begin{tabular}{@{} l l @{}}
\toprule
\textbf{Parameter} & \textbf{Value} \\
\midrule
Epochs              & 6 \\
Effective batch size & 16 (4$\times$4 gradient accumulation) \\
Optimizer           & AdamW \\
Learning rate       & $1 \times 10^{-3}$ with cosine annealing \\
Loss function       & Cross-entropy (unweighted) \\
Mixed precision     & Disabled (T4 GPU compatibility) \\
Device              & NVIDIA Tesla T4 (16\,GB) \\
Checkpoint strategy & Best Macro~F1 saved per model \\
\bottomrule
\end{tabular}
\end{table}


% =====================================================================
\section{Training Results}
\label{sec:training_results}

\subsection{Overall Performance}

Table~\ref{tab:training_results} summarizes the best validation metrics
achieved during training. Mamba achieves 1.55$\times$ higher Macro~F1 than
the second-best model (SSM) and 1.77$\times$ higher than the Transformer.

\begin{table}[ht]
\centering
\caption{Training-phase performance comparison across all four models.}
\label{tab:training_results}
\begin{tabular}{@{} l r r r r @{}}
\toprule
\textbf{Model} & \textbf{Best Macro F1} & \textbf{Final Train Acc}
               & \textbf{Final Val Loss} & \textbf{Training Time} \\
\midrule
\textbf{Mamba} & \textbf{42.91\%} & 55.26\% & 1.085 &  9.1\,hrs \\
SSM            & 27.61\%          & 52.42\% & 1.241 &  5.6\,hrs \\
Transformer    & 24.24\%          & 58.25\% & 1.206 & 0.25\,hrs \\
LSTM           & 18.80\%          & 60.19\% & 1.270 & 0.30\,hrs \\
\bottomrule
\end{tabular}
\end{table}

\begin{figure}[ht]
\centering
\includegraphics[width=0.75\textwidth]{\trainplot{fig3_model_comparison.png}}
\caption{Model comparison: Best Macro~F1 and peak memory across all four
architectures.}
\label{fig:model_comparison}
\end{figure}

\begin{figure}[ht]
\centering
\includegraphics[width=0.75\textwidth]{\trainplot{fig5_summary_table.png}}
\caption{Summary table of training metrics for all four models.}
\label{fig:summary_table}
\end{figure}


\subsection{Training Dynamics}

\subsubsection{Loss Convergence.}
Figure~\ref{fig:loss_curves} shows that Mamba exhibits the steepest loss
reduction (1.29$\rightarrow$1.07), indicating effective learning of
discriminative features. The Transformer's loss decreases minimally
(1.29$\rightarrow$1.20), suggesting it struggles with raw ECG waveform
modeling at this sequence length.

\begin{figure}[ht]
\centering
\includegraphics[width=0.75\textwidth]{\trainplot{fig1_loss_curves.png}}
\caption{Validation loss curves across 6~epochs for all four models. Mamba
shows the steepest convergence.}
\label{fig:loss_curves}
\end{figure}

\begin{figure}[ht]
\centering
\includegraphics[width=0.9\textwidth]{\trainplot{fig7_loss_per_model.png}}
\caption{Individual training and validation loss curves per model, showing
detailed convergence behavior.}
\label{fig:loss_per_model}
\end{figure}

\subsubsection{F1 Progression.}
Mamba's F1 rises sharply from 18.8\% to 42.9\% by epoch~5, then drops to
34.9\% at epoch~6 --- indicating the onset of overfitting
(Figure~\ref{fig:f1_curves}). The SSM shows gradual improvement
(18.8\%$\rightarrow$27.6\%), confirming that Mamba's selective mechanism is
critical for ECG feature extraction. LSTM's F1 remains flat at 18.8\% across
all epochs, indicating majority-class prediction.

\begin{figure}[ht]
\centering
\includegraphics[width=0.75\textwidth]{\trainplot{fig2_f1_curves.png}}
\caption{Macro~F1 progression during training. Mamba achieves its peak at
epoch~5 before overfitting.}
\label{fig:f1_curves}
\end{figure}

\subsubsection{Training Accuracy.}
Figure~\ref{fig:train_accuracy} shows the training accuracy curves. LSTM
achieves the highest training accuracy (60.2\%) despite the lowest F1 (18.8\%),
revealing that it has learned to predict the majority class (Normal) rather
than discriminating between all four categories.

\begin{figure}[ht]
\centering
\includegraphics[width=0.75\textwidth]{\trainplot{fig4_train_accuracy.png}}
\caption{Training accuracy curves. High LSTM accuracy is deceptive ---
the model predicts the majority class.}
\label{fig:train_accuracy}
\end{figure}


\subsection{Efficiency Analysis}

Figure~\ref{fig:efficiency} plots Macro~F1 against training time, revealing
the accuracy--cost trade-off. Mamba requires more training time
(9.1\,hrs vs.\ 0.25\,hrs for Transformer) due to its sequential for-loop
over 4,500 time steps. However, this is a \textit{training-time} limitation
--- during edge inference, Mamba maintains a recurrent state with $O(1)$
per-step computation.

\begin{figure}[ht]
\centering
\includegraphics[width=0.65\textwidth]{\trainplot{fig8_efficiency_scatter.png}}
\caption{Efficiency scatter: Macro~F1 vs.\ training time. Mamba achieves the
best accuracy at moderate computational cost.}
\label{fig:efficiency}
\end{figure}


% =====================================================================
\section{Inference Results}
\label{sec:inference}

To evaluate generalization, we run inference on the held-out test set of
1,706~ECG recordings using the best checkpoints from training.
Table~\ref{tab:inference_results} summarizes the results.

\begin{table}[ht]
\centering
\caption{Inference-time performance on the held-out test set.}
\label{tab:inference_results}
\begin{tabular}{@{} l r r r r @{}}
\toprule
\textbf{Model} & \textbf{Macro F1} & \textbf{Weighted F1}
               & \textbf{Accuracy} & \textbf{Inference Time} \\
\midrule
Transformer    & 18.59\% & 43.97\% & 59.38\% & 23.5\,s \\
LSTM           & 18.18\% & 37.05\% & 40.91\% & 10.6\,s \\
SSM            & 13.58\% & 13.46\% & 17.76\% & 78.5\,s \\
Mamba          & 11.73\% & 13.89\% & 29.25\% & 130.4\,s \\
\bottomrule
\end{tabular}
\end{table}

\begin{figure}[ht]
\centering
\includegraphics[width=0.75\textwidth]{\infplot{overall_metrics.png}}
\caption{Overall inference metrics comparison: Macro~F1, Weighted~F1, and
Accuracy for all four models on the test set.}
\label{fig:overall_metrics}
\end{figure}

\begin{figure}[ht]
\centering
\includegraphics[width=0.75\textwidth]{\infplot{summary_table.png}}
\caption{Summary table of all inference metrics.}
\label{fig:inf_summary}
\end{figure}


\subsection{Confusion Matrix Analysis}

Figure~\ref{fig:confusion} reveals the prediction patterns of each model:
\begin{itemize}
    \item \textbf{Transformer:} Predicts ``Normal'' for nearly all inputs
          (majority-class collapse), achieving high accuracy (59.4\%) that
          merely reflects the class prior.
    \item \textbf{Mamba:} Predicts ``Other'' predominantly, correctly
          classifying 486/491 Other samples but misclassifying 998/1010 Normal
          samples.
    \item \textbf{LSTM:} Shows mixed Normal/AFib predictions, suggesting it
          learns some cardiac rhythm features but cannot reliably discriminate.
    \item \textbf{SSM:} Similar to Mamba but with weaker class separation.
\end{itemize}

\begin{figure}[ht]
\centering
\includegraphics[width=0.85\textwidth]{\infplot{confusion_matrices (1).png}}
\caption{Confusion matrices for all four models. Each model collapses to
predicting 1--2 dominant classes.}
\label{fig:confusion}
\end{figure}


\subsection{Per-Class F1 Analysis}

Figure~\ref{fig:per_class_f1} shows the class-wise F1 breakdown:
\begin{itemize}
    \item \textbf{Normal:} Only Transformer achieves meaningful F1 (74.2\%) by
          predicting Normal by default.
    \item \textbf{Other:} Mamba achieves the highest F1 (44.5\%).
    \item \textbf{AFib:} Only LSTM achieves a non-zero F1 (11.5\%).
    \item \textbf{Noise:} All models achieve 0\% F1 --- with only 57 test
          samples, no model learned to detect noise.
\end{itemize}

\begin{figure}[ht]
\centering
\includegraphics[width=0.75\textwidth]{\infplot{per_class_f1.png}}
\caption{Per-class F1 scores. Severe class imbalance causes all models to
fail on minority classes.}
\label{fig:per_class_f1}
\end{figure}


\subsection{ROC and Precision--Recall Analysis}

ROC curves (Figure~\ref{fig:roc}) and precision--recall curves
(Figure~\ref{fig:pr}) provide further insight. AUC values near 0.5 for most
class--model combinations confirm random-like discrimination. Precision drops
sharply as recall increases, indicating that none of the models can maintain
precision while retrieving more positive samples from minority classes.

\begin{figure}[ht]
\centering
\includegraphics[width=0.85\textwidth]{\infplot{roc_curves.png}}
\caption{ROC curves (one-vs-rest) for all four models.  Curves near the
diagonal indicate near-random classification.}
\label{fig:roc}
\end{figure}

\begin{figure}[ht]
\centering
\includegraphics[width=0.85\textwidth]{\infplot{precision_recall.png}}
\caption{Precision--Recall curves for all four models. Low area under the
curve for minority classes confirms poor discriminative ability.}
\label{fig:pr}
\end{figure}


\subsection{Inference Speed}

Figure~\ref{fig:inf_speed} compares inference times on the T4~GPU. The
Transformer (23.5\,s) and LSTM (10.6\,s) benefit from optimized CUDA kernels
and parallelizable operations. Mamba (130.4\,s) is slowest due to its
sequential selective scan. \textbf{Note:} batch-mode GPU inference does not
reflect edge deployment, where Mamba's $O(1)$ recurrent state update per
sample would outperform the Transformer's $O(N^2)$ attention recomputation.

\begin{figure}[ht]
\centering
\includegraphics[width=0.6\textwidth]{\infplot{inference_speed.png}}
\caption{Inference time for 1,706 test samples on T4~GPU. Mamba's sequential
scan is slower in batch mode but advantageous in streaming edge deployment.}
\label{fig:inf_speed}
\end{figure}


% =====================================================================
\section{Analysis and Discussion}
\label{sec:analysis}

\subsection{Training vs.\ Inference Gap}

Table~\ref{tab:gap} highlights the significant gap between training-time
validation and test-time inference metrics.

\begin{table}[ht]
\centering
\caption{Gap between training validation and test-time inference Macro~F1.}
\label{tab:gap}
\begin{tabular}{@{} l r r r @{}}
\toprule
\textbf{Model} & \textbf{Train Val F1} & \textbf{Test F1} & \textbf{Gap} \\
\midrule
Mamba       & 42.91\% & 11.73\% & $-31.2$\,pp \\
Transformer & 24.24\% & 18.59\% & $-5.7$\,pp  \\
LSTM        & 18.80\% & 18.18\% & $-0.6$\,pp  \\
SSM         & 27.61\% & 13.58\% & $-14.0$\,pp \\
\bottomrule
\end{tabular}
\end{table}

This gap is attributed to three factors:
\begin{enumerate}
    \item \textbf{WeightedRandomSampler during training:} Training used
          oversampling of minority classes, creating artificially balanced
          batches. Inference runs on the true imbalanced distribution
          (59\%~Normal).
    \item \textbf{Limited epochs (6):} Mamba showed overfitting at epoch~5,
          while the Transformer barely learned. 20--50~epochs with early
          stopping would improve all models.
    \item \textbf{No class-weighted loss:} Unweighted cross-entropy fails to
          penalize minority-class misclassifications, allowing the gradient
          signal from majority classes to dominate.
\end{enumerate}


\subsection{Why Mamba Outperforms During Training}

The 15.3~percentage-point gap between Mamba (42.9\%) and baseline SSM
(27.6\%) isolates the effect of the \textit{selective mechanism}:

\begin{itemize}
    \item In Mamba, the discretization step~$\Delta$, matrices $B$ and $C$ are
          computed as functions of the input at each time step.
    \item A sharp QRS complex produces a large~$\Delta$, causing the model to
          ``remember'' the event; a flat baseline segment produces a
          small~$\Delta$, allowing the model to ``forget.''
    \item The fixed-coefficient SSM cannot distinguish between clinically
          relevant and irrelevant ECG segments.
\end{itemize}

\subsection{Why Transformers Struggle}

Three factors limit Transformer performance on raw ECG:
\begin{enumerate}
    \item \textbf{Positional encoding gap:} Sinusoidal embeddings were designed
          for discrete tokens, not continuous waveforms.
    \item \textbf{Attention dilution:} With $N{=}4{,}500$ time steps,
          self-attention is diluted across all positions. Clinically relevant
          QRS complexes span $\sim$15 time steps ($<0.4\%$ of the sequence).
    \item \textbf{No causal structure:} ECG analysis benefits from stateful,
          causal processing that accumulates evidence over time.
\end{enumerate}

\subsection{Why LSTM Fails}

LSTM's 18.8\% F1 equals the majority-class baseline, revealing: (i)~vanishing
gradients over 4,500~steps, (ii)~unidirectional processing prevents using
future context, and (iii)~64-dimensional hidden states may be insufficient for
the full variety of ECG morphologies.


% =====================================================================
\section{Edge Device Deployment Analysis}
\label{sec:edge}

\begin{table}[ht]
\centering
\caption{Inference complexity and edge deployment characteristics.}
\label{tab:edge}
\begin{tabular}{@{} l c c c c @{}}
\toprule
\textbf{Model} & \textbf{Per-step} & \textbf{Memory} &
\textbf{Parallel} & \textbf{Approx.\ Size} \\
\midrule
Mamba       & $O(1)$ & $O(d \times N)$ & No  & $\sim$200\,KB \\
Transformer & $O(N)$ & $O(N^2)$        & Yes & $\sim$180\,KB \\
LSTM        & $O(1)$ & $O(h)$          & No  & $\sim$140\,KB \\
SSM         & $O(1)$ & $O(d \times N)$ & No  & $\sim$192\,KB \\
\bottomrule
\end{tabular}
\end{table}

For continuous ECG monitoring at 300\,Hz, Mamba processes each new sample in
$O(1)$ by maintaining a recurrent state, whereas the Transformer must
recompute $O(N^2)$ attention over all previous samples. After 10~seconds
(3,000~samples), attention requires $\sim$9\,M operations per update. All
models fit comfortably on microcontrollers with 256\,KB+ RAM (e.g., ARM
Cortex-M4/M7, ESP32).


% =====================================================================
\section{Limitations}
\label{sec:limitations}

\begin{itemize}
    \item \textbf{Training duration:} Only 6~epochs due to Kaggle GPU
          constraints. Longer training with early stopping is expected to
          improve all models, particularly Mamba.
    \item \textbf{No hyperparameter tuning:} Identical hyperparameters were
          used for all models. Architecture-specific tuning may close the gap.
    \item \textbf{Class imbalance:} No focal loss, class weighting, or
          oversampling was applied. This disadvantages all models on minority
          classes.
    \item \textbf{Single-lead ECG:} CinC~2017 uses Lead~I only. Multi-lead
          ECG may favor architectures with more parameters.
    \item \textbf{Training speed:} Mamba's sequential scan makes it
          $37\times$ slower to train than the Transformer, though this does not
          affect per-sample inference on edge devices.
\end{itemize}


% =====================================================================
\section{Conclusion}
\label{sec:conclusion}

This study provides empirical evidence that Mamba's Selective State Space
architecture is better suited for ECG arrhythmia classification than
Transformers, LSTMs, and baseline SSMs. During training, Mamba achieves
42.91\%~Macro~F1 --- 55\% higher than the next-best model --- by leveraging
input-dependent state selection to dynamically focus on clinically relevant
ECG segments.

The gap between training and inference metrics is attributed to the training
pipeline (weighted sampling, limited epochs, unweighted loss) rather than the
architecture itself, and can be addressed through class-weighted loss, longer
training with early stopping, and proper validation protocols.

The $O(N)$ linear complexity, combined with $O(1)$ per-step recurrent
inference and a model size under 200\,KB, makes Mamba uniquely suited for
real-time ECG monitoring on edge devices where computational budgets are tight
and latency requirements are strict.

Future work should explore: (1)~longer training with class-weighted focal
loss, (2)~multi-lead ECG input, (3)~quantization and pruning for further edge
optimization, and (4)~hybrid Mamba--Transformer architectures.


% =====================================================================
\begin{thebibliography}{9}

\bibitem{gu2023mamba}
Gu, A., Dao, T.:
Mamba: Linear-Time Sequence Modeling with Selective State Spaces.
arXiv:2312.00752 (2023)

\bibitem{clifford2017af}
Clifford, G.D., et~al.:
AF Classification from a Short Single Lead ECG Recording: The
PhysioNet/Computing in Cardiology Challenge 2017.
Computing in Cardiology~\textbf{44} (2017)

\bibitem{vaswani2017attention}
Vaswani, A., et~al.:
Attention Is All You Need.
In: NeurIPS, vol.~30 (2017)

\bibitem{hochreiter1997lstm}
Hochreiter, S., Schmidhuber, J.:
Long Short-Term Memory.
Neural Computation~\textbf{9}(8), 1735--1780 (1997)

\bibitem{gu2022s4}
Gu, A., Goel, K., R{\'e}, C.:
Efficiently Modeling Long Sequences with Structured State Spaces.
In: ICLR (2022)

\end{thebibliography}

\end{document}
